\documentclass[12pt,letterpaper,reqno]{amsart}

\usepackage{epsfig}
\usepackage{amsmath}
\usepackage{amssymb}
\usepackage{amsthm}
\usepackage{indentfirst}
\usepackage{xspace}
\usepackage{multirow}
\usepackage{hyperref}
\usepackage{caption}
\usepackage{xcolor}
\usepackage{verbatim}
\usepackage[letterpaper,margin=1in,headheight=15pt]{geometry}
\usepackage{tikz-cd}
\usepackage{booktabs}
\usepackage{framed}
\usepackage{float}
\usepackage{thmtools}
\usepackage{dashrule}
\usepackage{silence}
\WarningFilter{gitinfo2}{I can't find}
\WarningFilter{biblatex}{Patching}
\usepackage[missing=]{gitinfo2}
\usepackage{fancyhdr}
\usepackage{mathtools}
\usepackage{enumitem}

\usepackage{libertine}
\renewcommand{\ttdefault}{lmtt}
\usepackage[nopatch=eqnum]{microtype}

\newcommand{\Cdotfill}{\leaders\hbox{$\,\cdot$}\hfill}
\newcommand{\aij}[1]{\tikzmarknode{a#1}{a_{#1}}}

\setlist[enumerate]{topsep=3.0pt,itemsep=0.4ex,partopsep=0.4ex,parsep=0.4ex}

\mathtoolsset{showonlyrefs}

\setlength{\footskip}{18pt}

\definecolor{purple}{rgb}{0.4,0.05,0.8}
\definecolor{darkbrown}{rgb}{0.5,0.5,0.5}
\definecolor{darkblue}{rgb}{0.1,0.1,0.6}
\definecolor{darkred}{rgb}{0.45,0.1,0.1}
\definecolor{darkgreen}{rgb}{0.0,0.3,0.06}
\hypersetup{colorlinks=true,urlcolor=darkred,linkcolor=purple,citecolor=darkred}
\definecolor{shadecolor}{rgb}{0.9,0.9,0.9}

% Bibliography formatting
\usepackage[bibstyle=numeric,defernumbers=true,maxnames=20,uniquename=init,backend=biber,sorting=none,bibencoding=utf8,casechanger=latex2e,backref=true]{biblatex}

%\DeclareNameAlias{sortname}{first-last}

\DeclareFieldFormat{url}{\url{#1}}
\DeclareFieldFormat[article]{pages}{#1}
\DeclareFieldFormat[inproceedings]{pages}{\lowercase{pp.}#1}
\DeclareFieldFormat[incollection]{pages}{\lowercase{pp.}#1}
\DeclareFieldFormat[article]{volume}{\textbf{#1}}
\DeclareFieldFormat[article]{number}{(#1)}
\DeclareFieldFormat[article]{title}{\MakeCapital{#1}}
\DeclareFieldFormat[inproceedings]{title}{#1}
\DeclareFieldFormat{shorthandwidth}{#1}

% Don't use "In:" in bibliography. Omit urls from journal articles.
\DeclareBibliographyDriver{article}{%
  \usebibmacro{bibindex}%
  \usebibmacro{begentry}%
  \usebibmacro{author/editor}%
  \setunit{\labelnamepunct}\newblock
  \MakeSentenceCase{\usebibmacro{title}}%
  \newunit
  \printlist{language}%
  \newunit\newblock
  \usebibmacro{byauthor}%
  \newunit\newblock
  \usebibmacro{byeditor+others}%
  \newunit\newblock
  \printfield{version}%
  \newunit\newblock
%  \usebibmacro{in:}%
  \usebibmacro{journal+issuetitle}%
  \newunit\newblock
  \printfield{note}%
  \setunit{\bibpagespunct}%
  \printfield{pages}
  \newunit\newblock
  \usebibmacro{eprint}
  \newunit\newblock
  \printfield{addendum}%
  \newunit\newblock
  \usebibmacro{pageref}%
  \usebibmacro{finentry}}

% Remove dot between volume and number in journal articles.
\renewbibmacro*{journal+issuetitle}{%
  \usebibmacro{journal}%
  \setunit*{\addspace}%
  \iffieldundef{series}
    {}
    {\newunit
     \printfield{series}%
     \setunit{\addspace}}%
  \printfield{volume}%
%  \setunit*{\adddot}%
  \printfield{number}%
  \setunit{\addcomma\space}%
  \printfield{eid}%
  \setunit{\addspace}%
  \usebibmacro{issue+date}%
  \newunit\newblock
  \usebibmacro{issue}%
  \newunit}

% Bibliography categories
\def\makebibcategory#1#2{\DeclareBibliographyCategory{#1}\defbibheading{#1}{\section*{#2}}}
\makebibcategory{books}{Books}
\makebibcategory{papers}{Refereed research papers}
\makebibcategory{chapters}{Book chapters}
\makebibcategory{conferences}{Papers in conference proceedings}
\makebibcategory{techreports}{Unpublished working papers}
\makebibcategory{bookreviews}{Book reviews}
\makebibcategory{editorials}{Editorials}
\makebibcategory{phd}{PhD thesis}
\makebibcategory{subpapers}{Submitted papers}
\makebibcategory{curpapers}{Current projects}

\setlength{\bibitemsep}{2.65pt}
\setlength{\bibhang}{.8cm}
\renewcommand{\bibfont}{\small}

\renewcommand*{\bibitem}{\addtocounter{papers}{1}\item \mbox{}\hskip-0.85cm\hbox to 0.85cm{\hfill\arabic{papers}.~~}}
\defbibenvironment{bibliography}
{\list{}
  {\setlength{\leftmargin}{\bibhang}%
   \setlength{\itemsep}{\bibitemsep}%
   \setlength{\parsep}{\bibparsep}}}
{\endlist}
{\bibitem}

\newenvironment{publications}{\section{\LARGE Publications}\label{papersstart}\vspace*{0.2cm}\small
\titlespacing{\section}{0pt}{1.5ex}{1ex}\itemsep=0.00cm
}{\label{papersend}\addtocounter{sumpapers}{-1}\refstepcounter{sumpapers}\label{sumpapers}}

\def\printbib#1{\printbibliography[category=#1,heading=#1]\lastref{sumpapers}}

% Counters for keeping track of papers
\newcounter{papers}\setcounter{papers}{0}
\newcounter{sumpapers}\setcounter{sumpapers}{0}
\def\lastref#1{\addtocounter{#1}{\value{papers}}\setcounter{papers}{0}}

% theorem environments
\declaretheoremstyle[spaceabove=0.25cm,spacebelow=0.25cm,notefont=\normalfont\bfseries, notebraces={(}{)}]{theorem}
\declaretheoremstyle[spaceabove=0.25cm,spacebelow=0.25cm,bodyfont=\normalfont,notefont=\normalfont\bfseries, notebraces={(}{)}]{noital}
\declaretheoremstyle[spaceabove=0.25cm,spacebelow=0.25cm,bodyfont=\normalfont\color{darkblue},notefont=\normalfont\bfseries, notebraces={(}{)}]{noitalblue}
\declaretheoremstyle[spaceabove=0.25cm,spacebelow=0.25cm,bodyfont=\normalfont\color{darkgreen},notefont=\normalfont\bfseries, notebraces={(}{)}]{green}
\declaretheoremstyle[spaceabove=0.25cm,spacebelow=0.25cm,bodyfont=\normalfont\color{darkred},notefont=\normalfont\bfseries, notebraces={(}{)}]{red}
\declaretheoremstyle[spaceabove=0.25cm,spacebelow=0.25cm,bodyfont=\normalfont\color{darkbrown},notefont=\normalfont\bfseries, notebraces={(}{)}]{noitaldarkbrown}
\declaretheoremstyle[spaceabove=0.25cm,spacebelow=0.25cm,bodyfont=\normalfont,notefont=\normalfont\bfseries,qed=$\qedsymbol$,notebraces={(}{)}]{proofstyle}

\newcounter{pset}

\declaretheorem[name=Theorem,numberwithin=section,style=red]{thm}
\declaretheorem[name=Conjecture,numberwithin=section,style=theorem]{conj}
\declaretheorem[name=Proposition,sibling=thm,style=red]{prop}
\declaretheorem[name=Proposition,sibling=thm,numbered=no]{prop*}
\declaretheorem[name=Construction,sibling=thm,style=theorem]{constr}
\declaretheorem[name=Corollary,sibling=thm,style=red]{cor}
\declaretheorem[name=Lemma,sibling=thm,style=red]{lem}
\declaretheorem[name=Definition,sibling=thm,style=noitalblue]{defn}
\declaretheorem[name=Definition-Proposition,sibling=thm,style=noitalblue]{defprop}
\declaretheorem[name=Example,sibling=thm,style=green]{example}
\declaretheorem[name=Exercise,numberwithin=section,style=noitaldarkbrown]{exercise}
% \declaretheorem[name=Proof,style=proofstyle,numbered=no]{pf}

\numberwithin{equation}{section}


% macros for convenience
\newcommand{\tops}{\texorpdfstring}

\newcommand{\nid}{\noindent}

\newcommand{\fa}{{\mathfrak a}}
\newcommand{\fp}{{\mathfrak p}}
\newcommand{\fk}{{\mathfrak k}}
\newcommand{\fg}{{\mathfrak g}}
\newcommand{\fh}{{\mathfrak h}}
\newcommand{\fn}{{\mathfrak n}}
\newcommand{\fq}{{\mathfrak q}}
\newcommand{\fm}{{\mathfrak m}}
\newcommand{\fr}{{\mathfrak r}}
\newcommand{\fu}{{\mathfrak u}}
\newcommand{\fG}{{\mathfrak G}}
\newcommand{\fgl}{{\mathfrak {gl}}}
\newcommand{\fsl}{{\mathfrak {sl}}}

\newcommand{\hfsl}{{\widehat\fsl}}
\newcommand{\hfg}{{\widehat\fg}}

\newcommand{\hC}{{\widehat C}}
\newcommand{\hS}{{\widehat S}}
\newcommand{\hX}{{\widehat X}}

\newcommand{\cC}{\ensuremath{\mathcal C}}
\newcommand{\cD}{\ensuremath{\mathcal D}}
\newcommand{\cW}{\ensuremath{\mathcal W}}
\newcommand{\cG}{\ensuremath{\mathcal G}}
\newcommand{\cB}{\ensuremath{\mathcal B}}
\newcommand{\cL}{\ensuremath{\mathcal L}}
\newcommand{\cS}{\ensuremath{\mathcal S}}
\newcommand{\cF}{\ensuremath{\mathcal F}}
\newcommand{\bcF}{\ensuremath{\overline{\mathcal F}}}
\newcommand{\cK}{\ensuremath{\mathcal K}}
\newcommand{\cZ}{\ensuremath{\mathcal Z}}
\newcommand{\cM}{\ensuremath{\mathcal M}}
\newcommand{\cP}{\ensuremath{\mathcal P}}
\newcommand{\cN}{\ensuremath{\mathcal N}}
\newcommand{\cO}{\ensuremath{\mathcal O}}
\newcommand{\cT}{\ensuremath{\mathcal T}}
\newcommand{\cH}{\ensuremath{\mathcal H}}
\newcommand{\cX}{\ensuremath{\mathcal X}}
\newcommand{\cY}{\ensuremath{\mathcal Y}}
\newcommand{\cA}{\ensuremath{\mathcal A}}
\newcommand{\cI}{\ensuremath{\mathcal I}}

\newcommand{\tphi}{{\ensuremath{\widetilde \Phi}}}
\newcommand{\tf}{{\ensuremath{\widetilde f}}}
\newcommand{\tg}{{\ensuremath{\widetilde g}}}
\newcommand{\tj}{{\ensuremath{\widetilde j}}}
\newcommand{\tiota}{{\ensuremath{\widetilde \iota}}}
\newcommand{\talpha}{{\ensuremath{\widetilde \alpha}}}
\newcommand{\tbeta}{{\ensuremath{\widetilde \beta}}}
\newcommand{\tU}{{\ensuremath{\widetilde U}}}
\newcommand{\tV}{{\ensuremath{\widetilde V}}}
\newcommand{\tM}{{\ensuremath{\widetilde M}}}
\newcommand{\tN}{{\ensuremath{\widetilde N}}}
\newcommand{\tomega}{{\ensuremath{\widetilde \omega}}}
\newcommand{\tF}{{\ensuremath{\widetilde F}}}

\newcommand{\Heis}{{\mathrm{Heis}}}
\newcommand{\Vir}{{\mathrm{Vir}}}

\newcommand{\bbR}{\ensuremath{\mathbb R}}
\newcommand{\bbC}{\ensuremath{\mathbb C}}
\newcommand{\bbP}{\ensuremath{\mathbb P}}
\newcommand{\bbZ}{\ensuremath{\mathbb Z}}
\newcommand{\bbQ}{\ensuremath{\mathbb Q}}
\newcommand{\bbA}{\ensuremath{\mathbb A}}
\newcommand{\bbT}{\ensuremath{\mathbb T}}
\newcommand{\bbH}{\ensuremath{\mathbb H}}
\newcommand{\bbI}{\ensuremath{\mathbb I}}
\newcommand{\bbS}{\ensuremath{\mathbb S}}
\newcommand{\bbF}{\ensuremath{\mathbb F}}

\newcommand{\bJ}{\ensuremath{\overline J}}
\newcommand{\bphi}{\ensuremath{\overline \phi}}

\newcommand{\regular}{{\mathrm{reg}}}

\newcommand{\half}{\ensuremath{\frac{1}{2}}}
\newcommand{\qtr}{\ensuremath{\frac{1}{4}}}
\newcommand{\bq}{{\mathbf q}}
\newcommand{\bbN}{{\mathbb N}}
\newcommand{\N}{{\mathcal N}}
\newcommand{\F}{{\mathcal F}}
\newcommand{\HH}{{\mathcal H}}
\newcommand{\LL}{{\mathcal L}}
\newcommand{\RR}{{\mathcal R}}
\newcommand{\V}{{\mathcal V}}
\newcommand{\dirac}{\!\!\not\!\partial}
\newcommand{\Dirac}{\!\!\not\!\!D}
\newcommand{\cE}{{\mathcal E}}
\newcommand{\vs}{\not\!v}
\newcommand{\kahler}{K\"ahler\xspace}
\newcommand{\painleve}{Painlev\'e\xspace}
\newcommand{\kq}{/\!\!/}
\newcommand{\kql}[1]{/\!\!/\!\!_#1\,}
\newcommand{\hk}{hyperk\"ahler\xspace}
\newcommand{\Hk}{Hyperk\"ahler\xspace}
\newcommand{\hkq}{/\!\!/\!\!/\!\!/}
\newcommand{\hkql}[1]{/\!\!/\!\!/\!\!/\!\!_#1\,}
\newcommand{\del}{\ensuremath{\partial}}
\newcommand{\delbar}{\ensuremath{\overline{\partial}}}
\newcommand{\vphi}{{\vec \Phi}}
\newcommand{\I}{{\mathrm i}}
\newcommand{\J}{{\mathrm j}}
\newcommand{\K}{{\mathrm k}}
\newcommand{\e}{{\mathrm e}}
\newcommand\bid{{\mathbf 1}}
\newcommand\triv{{\mathbf {triv}}}
\newcommand{\De}{\mathrm{D}}
\newcommand{\de}{\mathrm{d}}
\newcommand{\ab}{\mathrm{ab}}
\newcommand{\vol}{\mathrm{vol}}
\renewcommand{\sf}{\mathrm{sf}}
\newcommand{\inst}{\mathrm{inst}}
\newcommand{\eff}{\mathrm{eff}}
\newcommand{\fiber}{\mathrm{fiber}}
\newcommand{\rmvert}{\mathrm{vert}}
\newcommand{\rmcyc}{\mathrm{cyc}}
\newcommand{\closed}{\mathrm{closed}}
\newcommand{\exact}{\mathrm{exact}}
\newcommand{\ext}{\mathrm{ext}}
\newcommand{\op}{\mathrm{op}}
\newcommand{\Nek}{\mathrm{Nek}}
\newcommand{\oneloop}{\mathrm{1-loop}}
\newcommand{\p}{\partial}
\newcommand{\GL}{\mathrm{GL}}
\newcommand{\SL}{\mathrm{SL}}
\newcommand{\SU}{\mathrm{SU}}
\newcommand{\Mat}{\mathrm{Mat}}
\newcommand{\PGL}{\mathrm{PGL}}
\newcommand{\PSL}{\mathrm{PSL}}
\newcommand{\formal}{\mathrm{formal}}
\newcommand{\vertex}{\mathcal{V}}
\newcommand{\vrep}{\mathcal{R}}
\newcommand{\coinv}{{\mathrm{coinv}}}
\newcommand{\Coinv}{{\mathrm{Coinv}}}

\newcommand{\abs}[1]{\lvert#1\rvert}
\newcommand{\bigabs}[1]{\left\lvert#1\right\rvert}
\newcommand{\norm}[1]{\lVert#1\rVert}
\newcommand{\bignorm}[1]{\left\lVert#1\right\rVert}
\newcommand{\IP}[1]{\langle#1\rangle}
\newcommand{\DIP}[1]{\langle\!\langle#1\rangle\!\rangle}
\newcommand{\dwrt}[1]{\frac{\partial}{\partial#1}}
\newcommand{\eps}{\epsilon}
\newcommand{\simarrow}{\xrightarrow\sim}

\newcommand{\as}{{\text{ as }}}
\newcommand{\tin}{{\text{ in }}}

\newcommand{\Airy}{\textrm{Airy}}
\newcommand{\Betti}{\textrm{Betti}}
\newcommand{\Conf}{\textrm{Conf}}
\newcommand{\Fl}{\textrm{Fl}}
\newcommand{\Sol}{\textrm{Sol}}
\newcommand{\Mon}{\textrm{Mon}}
\newcommand{\dR}{\textrm{dR}}
\newcommand{\dRp}{\textrm{dR,p}}
\newcommand{\Ai}{\textrm{Ai}}
\newcommand{\Bi}{\textrm{Bi}}
\newcommand{\tC}{{\widetilde C}}
\newcommand{\tSigma}{{\widetilde \Sigma}}
\newcommand{\tS}{{\widetilde S}}
\newcommand{\near}{\textrm{near}}
\newcommand{\far}{\textrm{far}}
\newcommand{\std}{\textrm{std}}

\newcommand{\pd}[2]{\frac{\partial #1}{\partial #2}}

\newcommand{\mmaref}[1]{}

\newcommand{\ti}[1]{\textit{#1}}
\newcommand{\tb}[1]{\textbf{#1}}

\DeclareMathOperator{\sgn}{sgn}
\DeclareMathOperator{\supp}{supp}
\DeclareMathOperator{\ad}{ad}
\DeclareMathOperator{\im}{Im}
\DeclareMathOperator{\re}{Re}
\DeclareMathOperator{\Tr}{Tr}
\DeclareMathOperator{\Gr}{Gr}
\DeclareMathOperator{\Int}{Int}
\DeclareMathOperator{\Bd}{Bd}
\DeclareMathOperator{\End}{End}
\DeclareMathOperator{\Hom}{Hom}
\DeclareMathOperator{\Aut}{Aut}
\DeclareMathOperator{\Sym}{Sym}
\DeclareMathOperator{\Lie}{Lie}
\DeclareMathOperator{\diag}{diag}
\DeclareMathOperator{\Bun}{Bun}
\DeclareMathOperator{\Vect}{Vect}
\DeclareMathOperator{\Span}{Span}
\DeclareMathOperator{\grad}{grad}
\DeclareMathOperator{\rank}{rank}
\DeclareMathOperator{\ind}{ind}
\DeclareMathOperator{\coker}{coker}
\DeclareMathOperator{\Jac}{Jac}
\DeclareMathOperator{\Hol}{Hol}
\DeclareMathOperator{\Res}{Res}
\DeclareMathOperator{\gr}{gr}
\DeclareMathOperator{\Ob}{Ob}
\DeclareMathOperator{\Path}{Path}
\DeclareMathOperator{\LPath}{LPath}
\DeclareMathOperator{\id}{id}

\newcommand{\insfig}[3]{

\medskip
\noindent
\begin{minipage}{\linewidth}

\captionsetup{type=figure} 

\makebox[\linewidth]{\includegraphics[keepaspectratio=true,scale=#2]{figures/#1.pdf}}

\captionof{figure}{#3}
\label{fig:#1}

\end{minipage}
\medskip

}


% \newcommand{\insfig}[2]{\begin{figure}[htbp] \centering \includegraphics[scale=#2]{figures/#1-crop.pdf} \label{fig:#1} \end{figure}}
% syntax: \insfig{name}{0.5}{caption}

\DeclareRobustCommand{\fixme}[1]{\leavevmode{\color{orange}[#1]}}

% \mathtoolsset{showonlyrefs}

\bibliography{conformal-blocks}

\begin{document}

\makeatletter
\renewcommand{\sectionautorefname}{\S\@gobble}
\renewcommand{\subsectionautorefname}{\S\@gobble}
\renewcommand{\subsubsectionautorefname}{\S\@gobble}
\makeatother

\pagestyle{fancy}
\lhead{{\tiny \color{gray} \tt \gitAuthorIsoDate}}
\chead{\tiny \ti{Conformal blocks, spring 2026, Yale University}}
\rhead{{\tiny \color{gray} \tt \gitAbbrevHash}}
\renewcommand{\headrulewidth}{0.5pt}

\thispagestyle{fancy}

\begin{center}
\tb{Conformal blocks and some of their uses} \\
Andrew Neitzke
\end{center}

{These are notes from a spring 2026
course at Yale University, in progress.
\ti{Right now they are basically wreckage, but
should stabilize over the term.}

Please send corrections/improvements to
{\tt\href{mailto:andrew.neitzke@yale.edu}{andrew.neitzke@yale.edu}}.

% \tableofcontents

% \renewcommand{\listtheoremname}{Quick reference}
% \listoftheorems[onlynamed, numwidth=2.2em]

\thispagestyle{fancy}

\section{Introduction}

\subsection{Preface}

The course is about a subject which started in physics of
2-dimensional systems (partly string theory, but not only 
string theory). It has turned out to have a lot of applications,
in diverse areas.
\fixme{...}

\subsection{Vertex algebras and modules}

The first basic datum of the course is a 
\ti{vertex algebra} $\vertex$. 
A standard reference is
\cite{kac1998vertex}; our philosophy is close in spirit to
\cite{FrenkelBenZvi2004VertexAlgebrasCurves}.

A vertex algebra is a vector space $\vertex$ 
with additional structure.
We think of $a \in \vertex$ as objects which can be inserted at a point in the complex plane, giving an object $a(z)$ which varies analytically with 
$z$. 

$\vertex$ can act on a \ti{$\vertex$-module}: this means a vector space $\vrep$,
which again can be inserted at any point of the complex plane.
Let us put it at $z = 0$.
Then for any $a \in \vertex$, 
$a(z)$ should provide an endomorphism $Y^\vrep(a,z)$ of $\vrep$,
depending analytically on $z$.
We don't expect that $Y^\vrep(a,z)$ exists at $z = 0$; our motivation is
from local defects in quantum field theory, and often there are
singularities when two defects are placed at the same point.
But still we can expand in a series around $z = 0$.
Thus for $\chi \in \vrep$ we weill have
\begin{equation} \label{eq:y-series}
  Y^\vrep(a,z) \chi = \sum_{n \in \bbZ} z^{-n-1} a^\vrep_{(n)} \chi
\end{equation}
where each $a^\vrep_{(n)} \in \End \vrep$. 
One useful special case is that $\vertex$ is a module over itself
(``vacuum module''):
in that case we just write $Y(a,z)$ for $Y^\vertex(a,z)$,
and $a_{(n)}$ for $a^\vertex_{(n)}$.

Usually we consider $\vertex$-modules for which 
\eqref{eq:y-series} is a Laurent series,
i.e. $a^\vrep_{(n)} \chi = 0$ for $n \gg 0$, i.e. $Y^\vrep(a,z) \chi$ is formally meromorphic in $z$. In particular, we require
that $\vertex$ itself has this property.

When we have $a, b \in \vertex$ we also want to consider
a combined object $a(z) b(w)$. This 
object should again provide an endomorphism
of $\vrep$, depending
analytically on $z$ and $w$ away from $z = 0$ or $w = 0$, 
except that it can also have a pole at $z = w$.
This means that it may have different Laurent expansions
around $z = w = 0$ depending on whether
$\abs{w} < \abs{z}$ or $\abs{w} > \abs{z}$.\footnote{Recall the example
\begin{equation}
 f(z,w) = \frac{1}{z-w} = \begin{cases} z^{-1} \sum_{n \ge 0} (w/z)^n & \text{ for } \abs{w} < \abs{z}, \\ -w^{-1} \sum_{n \ge 0} (z/w)^n & \text{ for } \abs{w} > \abs{z}. \end{cases}
\end{equation}
}
We interpret these two expansions as $Y^\vrep(a,z) Y^\vrep(b,w) \chi$ and $Y^\vrep(b,w) Y^\vrep(a,z) \chi$ respectively.\footnote{
The idea is that the operator closer to $0$ acts first,
and then the operator farther from $0$ acts on the 
resulting object. We won't attempt to formalize this statement
precisely now; for physicists it is an instance of \ti{radial quantization}; it can be formalized in the language of \ti{holomorphic factorization algebras}, see e.g. \cite{CostelloGwilliam2021FA2}.}

Now we can make the definition:

\begin{defn}[Vertex algebra]

A \emph{vertex algebra} is a quadruple \((\vertex, \mathbf{1}, \partial, Y)\) where
\(\vertex\) is a vector space, \(\bid \in \vertex \),
\(T \in \mathrm{End}(\vertex)\),
and
\begin{equation}
Y(\,\cdot\,,z): \vertex \to \mathrm{End}(\vertex)[[z,z^{-1}]],\qquad
Y(a,z) = \sum_{n\in \mathbb{Z}} a_{(n)} z^{-n-1},
\end{equation}
is a linear map, such that:

\begin{enumerate}

\item For each \(a, b \in \vertex \) we have \(a_{(n)}b=0\) for all \(n\gg 0\) \, ,

\item  
$Y(\mathbf{1},z) = \bid_{\vertex}$ \, ,

\item
$Y(a,z)\mathbf{1}\in a+z\,\vertex[[z]]$ \, ,

\item
$[T,\,Y(a,z)] = \partial_z\, Y(a,z)$ \, ,

\item For all $a,b,\chi \in \vertex$
there exists $F(z,w) \in \vertex[[z,w]][z^{-1},w^{-1},(z-w)^{-1}]$, such that
$Y(a,z) Y(b,w) \chi$ and 
$Y(b,w) Y(a,z) \chi$ are the
formal Laurent expansion of $F(z,w)$ in the regions
$\abs{w} < \abs{z}$ and $\abs{z} < \abs{w}$
respectively.

\end{enumerate}

\end{defn}

% \begin{exercise}
% For all \(a,b\in\vertex\), prove that
% \begin{align}
% &z_0^{-1}\delta\!\left(\frac{z_1-z_2}{z_0}\right)Y(a,z_1)Y(b,z_2)
% - z_0^{-1}\delta\!\left(\frac{z_2-z_1}{-z_0}\right)Y(b,z_2)Y(a,z_1) \\
% &\qquad =
% z_2^{-1}\delta\!\left(\frac{z_1-z_0}{z_2}\right)Y\!\big(Y(a,z_0)b,\,z_2\big),
% \end{align}
% as an identity in \(\mathrm{End}(\vertex)\big[\!\big[z_0^{\pm1},z_1^{\pm1},z_2^{\pm1}\big]\!\big]\),
% where \(\delta(w)=\sum_{n\in\mathbb{Z}} w^n\) is the formal delta function.
% \end{exercise}

As you can perhaps see, a vertex algebra is a complicated object,
and a little painful to construct by hand. We will describe an
existence theorem below, which suffices for all our purposes.

The next few exercises develop some useful facts; I think you can
find them solved in \cite{FrenkelBenZvi2004VertexAlgebrasCurves}
Section 3.2.

\begin{exercise}
Prove the translation property
\begin{equation}
   \e^{w T} Y(a, z) \e^{-w T} = Y(a, z+w) \, .
\end{equation}
\end{exercise}

\begin{exercise}
Prove the skew-symmetry property
\begin{equation}
  Y(a,z) b = \e^{z T} Y(b, -z) a
\end{equation}
\end{exercise}

\begin{exercise}
Prove that the modes of $\vertex$ form a Lie subalgebra of $\End(\vertex)$, with
% determined by the singularity in the 
% products $a(z) b(w)$ at $z=w$, i.e. determined by the
% $a_{(n)} b$ for $n \ge 0$:
\begin{equation} \label{eq:mode-algebra}
[a_{(m)},\,b_{(n)}]
= \sum_{j\ge 0} \binom{m}{j}\,\bigl(a_{(j)}b\bigr)_{(m+n-j)} \, .
\end{equation}
% (Hint: represent the modes of $a$ and $b$ as contour integrals
% around $z=0$. Then calculate by residues at $z=w$. 
% Start with the simplest case where the only
% nonzero term is $j=0$.)
\end{exercise}
This is a useful way of thinking about vertex algebras
concretely, and particularly their ``local'' behavior: 
it is determined by the corresponding Lie algebra of modes.

\begin{defn}[Module for a vertex algebra]
A \ti{$\vertex$-module} is a vector space $\vrep$
and an operation
\begin{equation}
  Y^\vrep: \vertex \to \End \vrep[[z^{\pm 1}]]
\end{equation}
such that
\begin{enumerate}

\item 
  $Y^\vrep(\bid, z) = \bid_\vrep$ \, ,

\item $[T, Y^\vrep(a, z)] = \partial_z Y^\vrep(a,z)$ \, ,

\item For all $a,b \in \vertex$ and $\chi \in \vrep$
there exists $F(z,w) \in \vrep[[z,w]][z^{-1},w^{-1},(z-w)^{-1}]$, such that
$Y(a,z) Y(b,w) \chi$ and 
$Y(b,w) Y(a,z) \chi$ are the
formal Laurent expansion of $F(z,w)$ in the regions
$\abs{w} < \abs{z}$ and $\abs{z} < \abs{w}$
respectively.
\end{enumerate}

\end{defn}



% \begin{exercise} Prove that
% \begin{equation} \label{eq:symmetry-mode-products}
% a_{(n)}b = -\sum_{j\ge 0} (-1)^{n+j}\,\frac{\partial^{\,j}}{j!}\,\bigl(b_{(n+j)}a\bigr)
% \qquad (n\ge 0) \, .
% \end{equation}
% \end{exercise}

% \begin{exercise} The complicated shape of \eqref{eq:symmetry-mode-products} might convince you that
% there should be a more symmetric way of describing the
% singular part of the product. There is, and it's
% called the $\lambda$-bracket. Learn about it.
% \end{exercise}

% \begin{align}
% [a_\lambda b] &= \sum_{j\ge 0}\frac{\lambda^j}{j!}\,(a_{(j)}b),\\[4pt]
% a(z) &= \sum_{m\in\mathbb Z} a_{(m)}\,z^{-m-1},\\[4pt]
% \bigl[a_{(m)},\,b_{(n)}\bigr]
% &= \sum_{j\ge 0} \binom{m}{j}\,\bigl(a_{(j)}b\bigr)_{(m+n-j)}.
% \end{align}

% Heisenberg example:  [a_\lambda a]=\kappa\,\lambda  so a_{(1)}a=\kappa and a_{(j)}a=0 for j\neq 1.
% Hence
% If a_{(1)}a=\kappa\mathbf{1} (a constant), then (a_{(1)}a)_{(r)}=\kappa\,\delta_{r,-1},
% giving  [a_{(m)},a_{(n)}]= m\,\kappa\,\delta_{m+n,0}.


% getting a meromorphic expansion of the form 
% \begin{equation}
%  \psi(z) \phi(w) = \sum_{n \in \bbZ} 
%  (z-w)^{-n-1} (\psi_{(n)} \phi)(w) 
% \end{equation}
% where each $\psi_{(n)} \phi \in \vertex$.
% This operation is supposed to be commutative and 
% associative in a precise sense
% (which amounts to some complicated constraints on the products $(a,b) \mapsto a_{(n)} b$).
% The official definition is in \autoref{sec:vertex-algebras}.


\subsection{Universal enveloping vertex algebras}

Many of the vertex algebras we consider are of a simple sort:
they are \ti{universal enveloping vertex algebras},
which can be described efficiently by naming a finite number
of generating objects, and giving the singular part
of the expansion, i.e. giving the products $a_{(n)} b$ for 
$n \ge 0$.
We organize them into
the \ti{$\lambda$-bracket} defined by
\begin{equation}
  [a_\lambda b](w) = \Res_{z=w} \e^{\lambda(z-w)} a(z) b(w) = \sum_{j \ge 0} \frac{\lambda^j}{j!} a_{(j)} b \, .
\end{equation}
This gives the following structure on the generating fields:

\begin{defn}[Lie conformal algebra]
A \ti{Lie conformal algebra} is a $\bbC[\partial]$-module $R$
with a $\lambda$-bracket
\begin{equation}
  [\cdot_\lambda\cdot]: R \otimes R \to R[\lambda] \, ,
\end{equation}
expanded as
\begin{equation}
  [a_\lambda b] = \sum_{j \ge 0} \frac{\lambda^j}{j!} a_{(j)} b \, ,
\end{equation}
such that
\begin{gather}
  [\partial a_\lambda b] = -\lambda[a_\lambda b], \quad [a_\lambda \partial b] = (\lambda + \partial)[a_\lambda b], \quad [a_\lambda b] = -[b_{-\lambda-\partial} a] \, , \\
  [a_\lambda[b_\mu c]] - [b_\mu[a_\lambda c]] = [[a_\lambda b]_{\lambda+\mu} c] \, .
\end{gather}
\end{defn}

Any vertex algebra is a Lie conformal algebra, and we can
define a ``free'' extension of a given Lie conformal algebra
to a vertex algebra:

\begin{prop}[Existence of universal enveloping vertex algebra]
Suppose $R$ is a Lie conformal algebra. Then there exists
a vertex algebra $\vertex(R)$, such that for any 
vertex algebra $\vertex'$ and any morphism
of Lie conformal algebras $\phi: R \to \vertex'$,
there is a unique morphism of vertex algebras
$\Phi: \vertex(R) \to \vertex'$, such that for all $a \in R$,
\begin{equation}
  \Phi(a_{(-1)} \bid) = \phi(a) \, .
\end{equation}
\end{prop}

Here are three main examples: 
\begin{enumerate}
\item The \ti{Heisenberg algebra} $\vertex = \Heis$. It is generated by one defect $J$, i.e. $R = \bbC \cdot J$, 
and has operator product of the form
\begin{equation}
 J(z) J(w) = \frac{\bid(w)}{(z-w)^2} + \regular
\end{equation}
i.e. $[J_\lambda J] = \lambda \bid$.
% i.e. the only nonzero product is $J_{(1)} J = \bid$.
The mode algebra is spanned by $J_{(n)}$ $(n \in \bbZ)$ and $\bid = \bid_{(-1)}$, with
\begin{equation}
[J_{(m)},J_{(n)}] = m \delta_{m+n,0} \bid \, .
\end{equation}
% Because $Y(J,z) \bid \in J + z \vertex[[z]]$
% we have the relations $J_{(n)} \bid = 0$ for $n \ge 0$.

% As a module over the Lie algebra of modes, $\vertex$
% is freely generated by $\bid$ modulo these relations. Thus
% a basis for $\Heis$ is provided by elements
% $J_{(n_1)} \cdots J_{(n_k)} \bid$ with $n_1 \le n_2 \le \cdots \le n_k < 0$.

\item For any complex Lie algebra $\fg$ with invariant bilinear form $\kappa$ and $k \in \bbC$, 
one has the \ti{affine Kac-Moody algebra} $\vertex = V^k(\fg)$.
It is generated by defects $J_X$ for every $X \in \fg$, i.e.
$R = \fg$, with operator products of the form
\begin{equation} \label{eq:kac-moody-ope}
  J^X(z) J^Y(w) = \frac{k \kappa(X,Y) \bid(w)}{(z-w)^2} + \frac{J^{[X,Y]}(w)}{z-w} + \regular
\end{equation}
i.e. $[J^X\,\!_\lambda J^Y] = \lambda k \kappa(X,Y) \bid + J^{[X,Y]}$.
The corresponding mode algebra is
\begin{equation}
[J^X_{(m)},J^Y_{(n)}] = J^{[X,Y]}_{(m+n)} + m k \kappa(X,Y) \delta_{m+n,0} \bid \, .
\end{equation}
Note the Heisenberg algebra is a special case, where $\fg$ is
abelian of rank $1$.

\item For any $c \in \bbC$ one has the \ti{Virasoro algebra} $\Vir^c$. It is generated by a defect $T$, with an expansion of the form
\begin{equation} \label{eq:virasoro-ope}
  T(z) T(w) = \frac{\frac{c}{2} \bid(w)}{(z-w)^4} + \frac{2 T(w)}{(z-w)^2} + \frac{\partial T(w)}{z-w} + \regular
\end{equation}
i.e. $[T_\lambda T] = \lambda^3 \frac{c}{12} \bid + (2 \lambda + \partial) T$.
The corresponding mode Lie algebra is
\begin{equation} \label{eq:virasoro-mode-algebra}
[L_{(m)},L_{(n)}]
=(m-n)\,L_{(m+n-1)}+\frac{c}{12}\,m(m-1)(m-2)\,\delta_{m+n,2} \bid \, .
\end{equation}
\end{enumerate}

\begin{exercise}
Derive the Lie algebra \eqref{eq:virasoro-mode-algebra}
from \eqref{eq:virasoro-ope} and \eqref{eq:mode-algebra}.
\end{exercise}


\subsection{Riemann surfaces}

The second main datum of the course is a \ti{Riemann surface} $C$, 
i.e. a complex manifold of complex dimension $1$. 
We usually take $C$ to be compact. 
Two fundamental examples are 
\begin{enumerate}
\item the \ti{Riemann sphere} $C = \bbC \bbP^1$, covered by two patches with coordinates $z$ and $z'$, related by $z' = 1/z$.
\item the \ti{torus} $C = \bbC / (\bbZ \oplus \tau \bbZ)$, depending
on the parameter $\tau \in \bbC$ with $\im \tau > 0$.
\end{enumerate}

We also often include marked points $p_i$ on $C$, each 
equipped with a local coordinate around it. 
Each marked point $p_i$ will be labeled by a $\vertex$-module 
$\vrep_i$.

We often drop the data of the marked points and the representations from the notation, just writing the whole object as $C$.


% \subsection{Conformal blocks}

% Given these data there is a space of \ti{conformal blocks}
% \begin{equation}
%   \Conf(C, \vertex) \, .
% \end{equation}
% The spaces $\Conf(C, \vertex)$ are the main object of study in this course. In this section I'll tell you some things about them,
% but not give a fully general definition yet.


\subsection{Chiral correlation functions for affine
Kac-Moody algebras}

Now we are ready to ``put $\vertex$ on the Riemann surface $C$.''
The way this is done for a general $\vertex$ takes a while
to appreciate. To get some examples in our hand quickly, 
let me specialize to the case of affine Kac-Moody algebras,
and further to the case of a curve $C$ with no marked points.

Before defining conformal blocks we consider correlation functions.
I think that these are actually the more important objects (and in good
cases they are the same thing as conformal blocks).

\begin{defn}[Chiral correlation system for $V^k(\fg)$, without marked points]
Suppose given a Riemann surface $C$.
A \ti{chiral correlation system} for $V^k(\fg)$
on $C$ is a rule that assigns an element
\begin{equation}
  \bigg\langle \prod_{j=1}^m J^{X_j}(q_j) \bigg\rangle \in \bigotimes_{j=1}^m T^*_{q_j} C  
\end{equation}
for every $m \ge 0$, $q_j \in C$,
which:
\begin{itemize}
  \item is linear in each $X_i$ and in each $\chi_j$,
  \item is symmetric under the interchange of the $J^{X_i}(q_i)$,
  \item is meromorphic in the $q_i$, and holomorphic away from
  loci where some $q_i = q_j$,
%   \item has Laurent expansion around $q_j = p_i$ given by
% \begin{equation}
%   \IP{J^X(q) \chi(p) \cdots} = \sum_{n \in \bbZ} z^{-n-1} \IP{(J^X_{(n)} \chi)(p) \cdots} \, ,
% \end{equation}
  \item satisfies the relations \eqref{eq:kac-moody-ope} around $q_i = q_j$, i.e.
  the polar part of
$\IP{J^X(q) J^Y(q') \cdots}$ around $q=q'$ is 
\begin{equation}
  \IP{J^X(q) J^Y(q') \cdots} = \frac{\IP{\cdots} \de z(q) \de z(q')}{(z(q)-z(q'))^2} + \frac{\IP{J^{[X,Y]}(q) \cdots} \de z(q')}{z(q)-z(q')} + \cdots \, .
\end{equation}
\end{itemize}

\end{defn}

% As we defined it, this is a highly complicated and structured object. We will see shortly 
% that some of the structure is redundant: the whole correlation
% system is determined by a relatively small subset of the
% correlation functions.
% But in the applications of the theory it is important to have
% all the correlation functions available.

\begin{example}[Chiral correlation systems for the Heisenberg vertex algebra on the sphere]

Here is a baby example.

Suppose $\vertex = \Heis = V^1(\fgl_1)$, and $C = \bbC \bbP^1$. 
Then we can construct a chiral correlation system as follows.
Let $\sigma$ denote a division of $\{1, \dots, k\}$ into 
2-element subsets $(\sigma_{i1}, \sigma_{i2})$ with $i = 1, \dots, \frac{k}{2}$. Then we can determine a chiral correlation system by
\begin{equation}
  \IP{J(z_1) \cdots J(z_k)} = \sum_\sigma \prod_{i=1}^k \frac{\de z_{\sigma_{i1}} \de z_{\sigma_{i2}}}{(z_{\sigma_{i1}} - z_{\sigma_{i2}})^2}
\end{equation}
so e.g.
\begin{align}
 \IP{} &= 1, \\
 \IP{J(z_1)} &= 0, \\
 \IP{J(z_1) J(z_2)} &= \frac{1}{(z_1 - z_2)^2} \de z_1 \de z_2, \\
 \IP{J(z_1) J(z_2) J(z_3)} &= 0, \\
 \IP{J(z_1) J(z_2) J(z_3) J(z_4)} &= \left( \frac{1}{(z_1 - z_2)^2 (z_3 - z_4)^2} + \frac{1}{(z_1 - z_3)^2 (z_2 - z_4)^2} + \frac{1}{(z_1 - z_4)^2 (z_2 - z_3)^2}  \right)  \de z_1 \de z_2 \de z_3 \de z_4, 
\end{align}


\fixme{...}
\end{example}

\begin{exercise} Prove that every 
chiral correlation system for this $(C,\vertex)$
is a multiple of the given one.
\end{exercise}

\begin{example}[Chiral correlation systems for the Heisenberg vertex algebra on the torus]
\fixme{...}
\end{example}



\subsection{Invariant spaces}

Suppose $\vertex = V^k(\fg)$, for any $k \in \bbC$.
Then for each finite-dimensional representation $R$ of $\fg$,
there is a corresponding infinite-dimensional representation $M(R)$ of $V^k(\fg)$.

Now take $C = \bbC\bbP^1$, with marked points $z_1, \dots, z_n$, carrying representations
$M(R_1), \dots, M(R_n)$. 
Then one has\footnote{The isomorphism we get here depends on $z_1 ,\dots, z_n$;
we'll see this from another point of view momentarily. It also depends on $k$.} \footnote{Proof idea: we can find a function $f$ realizing any polar parts; use their $X(f)$ to kill all the 
descendants, leaving
only the $R_\lambda$; what remains is the action of $X(1)$
on $\bigotimes_i R_i$.}

\begin{prop}[Conformal blocks for \tops{$V^k(\fg)$}{Vk(g)} give tensor product invariants]
\begin{equation}
\Conf(C, \vertex) = \Hom(\bigotimes_i R_i, \triv).
\end{equation}
\end{prop}

With 3 punctures this is equivalent to the tensor product multiplicity spaces, i.e.
\begin{equation}
\Conf(C, \vertex) = \Hom(R_1 \otimes R_2, R_3^*).
\end{equation}
But $\Conf(C,\vertex)$ has lots of extra structures, so we get some new
structure in representation theory this way.


\subsection{Connections}

Suppose that $\vertex$ is \ti{conformal}, i.e. it 
contains a copy of $\Vir^c$. For instance, this is 
true for $V^k(\fg)$ as long as $k \neq -h^\vee$.

Then we can consider varying the base curve $C$, and the spaces $\Conf(C, \vertex)$
will turn out to be the fibers of a twisted $D$-module $\Conf(\cdot, \vertex)$ over some moduli space $\cM$ of curves.
(Here we include the marked points as part of the data of $C$ --- so a special case of varying $C$
is to hold the bare curve fixed and just move the marked points.)
Roughly one should think of $\Conf(\cdot, \vertex)$ as a bundle with connection $\nabla$ over $\cM$.

Thus we get in particular some new structure on the
invariant spaces above \fixme{...}


\subsection{Link invariants}

In fact you get a little more, if $\vertex$ is sufficiently nice, e.g. \ti{strongly rational}. 

To get an example of a strongly rational vertex algebra
consider $V^k(\fg)$; it is simple for generic $k$, meaning
that as a module over itself it contains no proper submodule; 
for $k \in \bbZ_+$ it is not simple\footnote{A simple example: in $V^k(\fsl_2)$ the element $E_{-1}^k \bid$ generates a proper
submodule.} 
but it has a simple quotient $L_k(\fg)$, and $L_k(\fg)$ is strongly rational. Similarly $\Vir^c$ has strongly rational quotients
for special values of $c$.

If $\vertex$ is strongly rational, then it has only finitely many irreducible representations.\footnote{Fine print: I really mean \ti{ordinary modules}.}
Moreover the spaces $\Conf(C, \vertex)$ are finite-dimensional, 
for any compact $C$, with any irreducible representations inserted at punctures. So this is a nice world where all vector spaces
around are finite-dimensional, and we do not have to worry
about functional analysis, double duals, ind- or pro-objects, etc.

Now suppose $C$ is $\bbC\bbP^1$ with $n$ 
marked points $z_1, \dots, z_n$ in the finite plane, each one carrying a representation of $\vertex$. 
Then consider
moving any marked point $z_i$ along any loop in $C \setminus \{z_1, \dots, \hat{z}_i, \dots, z_n \}$.
The parallel transport of the connection $\nabla$ along 
this loop then gives an automorphism of $\Conf(C,\vertex)$.
Thus $\Conf(C,\vertex)$ carries an action of the braid group on $n$ strands.

Moreover you have, in addition to the braiding, ``cup'' and ``cap''
operations which create/destroy a pair of marked points.
Combining these you can build 
framed link invariants: indeed, letting $C$ be
$\bbC\bbP^1$ with \ti{no} marked points, each link in $\bbR^3$
gives a linear operator acting on a finite-dimensional space $\Conf(C, \vertex)$, 
and its trace is a link invariant.

One basic example: if $\vertex = L_k(\fsl_2)$ with $k \in \bbZ_+$, then the link invariant we get
is the \ti{Jones polynomial} $J_L(q)$, with $q = \exp(2 \pi \I / (k+2))$.

In this case
the representations $\vrep$ are objects of a \ti{modular tensor category}.


\subsection{Holomorphic factorization} 

In a 2-dimensional conformal field theory
one computes correlation functions of observables. 
These correlation functions are almost factorized into products of ``holomorphic'' and ``antiholomorphic'' parts,
but not exactly factorized; the nontrivial space of conformal blocks is the obstruction to factorization.

For example, in the CFT of the \ti{free boson} $\Phi$,
one is interested in correlation functions of the form
\begin{equation}
  \DIP{J(z_1) \cdots J(z_n) \bJ(w_1) \cdots \bJ(w_m)} 
\end{equation}
where classically $J = \partial \Phi$ and $\bJ = \overline\partial \Phi$.
Then one finds an approximate factorization
\begin{equation}
  \DIP{J(z_1) \cdots J(z_n) \bJ(w_1) \cdots \bJ(w_m)} = \int \de a \, \IP{J(z_1) \cdots J(z_n)}_{\Psi_a} \IP{\bJ(w_1) \cdots \bJ(w_m)}_{\overline\Psi_a}
\end{equation}
where $\IP{J(z_1) \cdots J(z_n)}_{\Psi_a}$ are holomorphic in a strong sense --- they
are holomorphic in the $z_i$, in $a$, and in the moduli of $C$.
They arise as part of the data of a conformal block $\Psi_a \in \Conf(C, \Heis)$.
So the conformal blocks $\Psi_a$ are the basic building-blocks one uses to build the 
correlation functions of the CFT.

For the simplest example of this, suppose $C$ is the flat torus with modulus 
$\tau$, $C = \bbC / (\bbZ \oplus \bbZ \tau)$, and any area $A$.
The simplest correlation function of the free boson is the ``$0$-point function'' $\DIP{1}$, which turns out to be given by
(after zeta-function regularization and deleting the zero mode,
see e.g. \cite{McIntyreTakhtajan2006})\footnote{In a CFT of central charge $c$, the partition function on a surface of Euler characteristic $\chi$ scales 
as $A^{\frac{c}{12} \chi}$ under overall rescaling of the metric. Here $\chi = 0$ so we should have no $A$ dependence in the end;
but $\det \Delta'$ 
scales as $A$, because $\zeta_\Delta(0) = \frac{1}{6} \chi - 1$,
the $-1$ coming from the zero mode;
this requirement tells us that we should compensate the removal
of the zero mode by inserting the factor $\sqrt{A}$ in the 
numerator. 
We can probably get the same factor from proper treatment of
the shift symmetry $\Phi \to \Phi + \alpha$? 
}


% BRST treatment
% of the shift symmetry?
% Write $Q \Phi = c$, $Q \bar{c} = b$, $Q c = Q b = 0$,
% $G(\Phi) = \int_C \star \Phi$, so 
% $S_{GF} = Q (\bar{c} G(\Phi)) = b \int_C \star \Phi - A \overline{c} c$,


% Or, slightly more abstractly: 
% we attempt to define the path integral using the $L^2$ measure
% over all modes. Then the integral over nonzero modes produces
% the factor $1 / \sqrt{\det \Delta'}$, and we still have
% remaining the $L^2$ measure on the zero modes.
% So the result of the path integral is best understood as
% $\frac{\mu_{L^2}}{\sqrt{\det \Delta'}}$. The zero mode field 
% $\phi = 1 / \sqrt{A}$ is properly normalized, so
% if we write $\phi = x_0 / \sqrt{A} + \cdots$ then 
% $\mu_{L^2} = \de x_0$. The total integral of such $\phi$
% against the area measure is $y_0 = \sqrt{A} x_0$,
% so $\mu_{L^2} = \sqrt{A} \, \de y_0$.
% Thus the result of the path integral is
% $\frac{\sqrt{A}}{\sqrt{\det \Delta'}} \de y_0$,
% and $\DIP{1}$ is scale invariant when measured
% in units of $\de y_0$.

\begin{equation} \label{eq:free-boson-partition-function-torus}
  \DIP{1} = \frac{\sqrt{A}}{\sqrt{\det \Delta'}} = \frac{1}{\sqrt{\im \tau} \abs{\eta(\tau)}^2}
  \end{equation}
Here $\eta(\tau)$ is the Dedekind eta-function,
\begin{equation}
  \eta(\tau) = q^{1/24} \prod_{n \ge 1} (1 - q^n), \qquad q = \e^{2 \pi \I \tau} \, .
\end{equation}
The answer \eqref{eq:free-boson-partition-function-torus} is \ti{almost} factorized into holomorphic and antiholomorphic, 
but not quite, because
of the factor $\sqrt{\im \tau}$.
The explanation of this in terms of conformal blocks is: the Heisenberg algebra on $C$
has conformal blocks $\Psi_a$ parameterized by $a \in \bbC$, with the correlation function
\begin{equation}
  \IP{1}_{\Psi_a} = \frac{1}{\eta(\tau)} \exp\left(\frac{1}{4 \pi \I} a^2 \tau\right)
\end{equation}
The blocks relevant for the free boson are the ones with $a \in \I \bbR$.
Then we have
\begin{equation} \label{eq:factorization-free-boson}
  \DIP{1} = \int_{a \in \I \bbR} \left(  \frac{\de a}{2 \pi \I} \right) \, \abs{\IP{1}_{\Psi_a}}^2  = \frac{1}{\sqrt{2 \im \tau} \abs{\eta(\tau)}^2}
\end{equation}
There is a similar, but more elaborate, story for 
higher genus surfaces \cite{McIntyreTakhtajan2006}.


\subsection{Liouville theory}

A more sophisticated example is provided by another 2d
CFT, \ti{Liouville theory}. In this case the relevant
vertex algebra is $\Vir^c$ rather than $\Heis$.

One needs a bit of notation: introduce a parameter $b$
related to the central charge by
\begin{equation}
c = 1 + 6 Q^2, \qquad Q = b + \frac{1}{b}
\end{equation}
and also let
\begin{equation}
  \Delta = j(Q-j) \, , j = \frac12 Q + \I p
\end{equation}

Now, for a simple example,
consider $C_q = \bbC \bbP^1$ with $4$ punctures at 
$0,q,1,\infty$, with Verma modules placed at the punctures.
Then there exists a family of conformal blocks
$\Psi_p \in \Conf(C_q, \Vir^c)$ \fixme{...}
for which we have a factorization formula similar to
\eqref{eq:factorization-free-boson}:\footnote{This is not the most common way of writing the factorization formula, but it can be found e.g. in \cite{Jackson:2014nla}.}
\begin{equation}
  \DIP{1} = \int_{p \in \bbR} (\rho(p) \, \de p) \, \abs{\IP{1}_{\Psi_p}}^2
\end{equation}
where we have the ``Plancherel measure''
\begin{equation}
  \rho(p) = 4 \sinh(2 \pi b p) \sinh(2 \pi b^{-1} p) \, ,
\end{equation}
and similar formulas for correlation functions.
\fixme{try to be more explicit, write DOZZ etc, understand better the issues about absorbing structure constant}


\subsection{Tau functions}

Let $\vertex$ be the Virasoro algebra with central charge $c=1$, $\vertex = \Vir_{c=1}$.
Then $\Conf(C, \vertex)$ is a module over a commutative algebra $\cA$, the algebra of
\ti{Verlinde loop operators}.
The joint eigenvectors of $\cA$ are particularly interesting.

For instance, let $C_q$ be $\bbC\bbP^1$ with $4$ punctures
placed at points $(z_1, z_2, z_3, z_4) = (0,1,\infty,q)$. At $z_i$ we place a simple module for $\Vir_{c=1}$:
a Verma module with highest weight $h_i$.
Then consider a 1-parameter family of conformal blocks $\Psi(q) \in \Conf(C_q, \Vir_{c=1})$
which are parallel, $\nabla_{\partial_q} \Psi(q) = 0$. 
Suppose moreover that $\Psi(q)$ is a joint eigenvector of $\cA$ for some $q$;
then it is a joint eigenvector for all $q$.\footnote{The algebra $\cA$ depends on the surface $C$, so here it should
really be written $\cA_q$ rather than just $\cA$; but the dependence is locally constant as we vary the surface.}
Let $\tau(q)$ be the $0$-point function
\begin{equation}
\tau(q) = \IP{1}_{\Psi(q)}
\end{equation}
Then $\tau(q)$ is a $\tau$-function of Painleve VI.
This means it is a solution of an interesting differential equation \fixme{...}




\subsection{The AGT correspondence} \fixme{...}

The AGT correspondence gives a very different interpretation of
the conformal blocks which showed up in Liouville theory.
We describe it in the concrete case of the 4-punctured
sphere $C_q$ which appeared above. Then
\begin{equation}
  \IP{1}_{\Psi_a} = Z_\Nek(a,q,m,\eps_1,\eps_2)
\end{equation}
The right side is the \ti{Nekrasov instanton partition function}
for $\cN=2$ supersymmetric Yang-Mills theory, with
gauge group $G = \SU(2)$, and matter content given by
$4$ copies of the fundamental representation:
\begin{equation}
  Z_\Nek = Z_{\oneloop} Z_\inst
\end{equation}
Here $Z_{\oneloop}$ is a $q$-independent function
\begin{equation}
Z_{\oneloop}
= \exp \Big(
\gamma_{{\eps_1},{\eps_2}}(2a;\Lambda)+\gamma_{{\eps_1},{\eps_2}}(-2a;\Lambda)
+\sum_{f=1}^4 \gamma_{{\eps_1},{\eps_2}}(a+m_f;\Lambda)+\gamma_{{\eps_1},{\eps_2}}(-a+m_f;\Lambda) \Big) \, ,
\end{equation}
in terms of a special function $\gamma_{\eps_1,\eps_2}(x;\Lambda)$
defined in \fixme{...},
and $Z_{\inst}$ is an explicit series in powers of $q$,
whose coefficients are rational functions of 
$a, \eps_1, \eps_2, m_f$. The surprising thing is that the
series $Z_{\inst}$ also arises in a totally different way,
as a generating series of integrals over the moduli space of
$\SU(2)$-instantons in $\bbR^4$.

The dictionary is
\begin{equation}
b = \frac{\eps_1}{\eps_2}
\end{equation}



\fixme{give explicit formula}


\subsubsection{Random matrix models} \fixme{...}

% Various random matrix models are also related to conformal blocks.
% Here we just discuss the simplest case.
% We consider the Hermitian $N \times N$ 
% $1$-matrix model: this means 

\subsubsection{Quantum groups} \fixme{...}

\subsubsection{Geometric Langlands} \fixme{...}

\subsubsection{Quantization of Teichm\"uller space} \fixme{...}

\subsubsection{Abelianization} \fixme{...}


% \section{Definition of conformal blocks}

% \fixme{...}


\section{The free scalar field CFT}

\subsection{The action and its conformal invariance}

Let $C$ be a Riemannian 2-manifold, with metric $g$.

We begin with a basic example of a Euclidean, 2-dimensional conformal field theory:
the (uncompactified) \ti{free scalar field}.
It has been made rigorous under the name
\ti{Gaussian multiplicative chaos}.
See e.g. \cite{GuillarmouRhodesVargas2019Polyakov2d}.
But we will stick to the more informal treatment.

We consider the (loosely understood) space of fields on $C$, 
\begin{equation}
  \cF = \left\{ \Phi: C \to \bbR \, \Bigg\lvert  \int_C \Phi \, \de \vol = 0 \right\} \, .
\end{equation}
We are interested in studying ``expectation values'' of random observables on $\cF$, formally
\begin{equation}
  \IP{\cO} = \int_{\Phi \in \cF} \de \mu(\Phi) \, \cO(\Phi) \, \e^{-S(\Phi)}
\end{equation}
where $\mu$ is a measure on $\cF$,
 $\cO$ is some function on $\cF$, and the \ti{action} $S: \cF \to \bbR$ is\footnote{We always use the Einstein summation convention of summing over repeated indices; cotangent indices are written as 
 subscripts and tangent indices are written as superscripts; thus $g_{ij}$ is the metric on $C$ in components i.e. $g_{ij} = g(\partial_{x^i}, \partial_{x^j})$, 
 and $g^{ij}$ is the inverse metric in components.}
\begin{equation} \label{eq:free-boson-action}
  S(\Phi) = \frac{1}{8\pi} \int_C \norm{\de \Phi}^2 \, \de \vol 
  % = \frac{1}{8 \pi} \int_C \de \Phi \wedge \star \de \Phi 
  % = -\frac{1}{8\pi}\int_C \Phi \, \de \star \de \Phi 
  = \frac{1}{8\pi}\int_C \partial_i \Phi \, \partial_j \Phi \, g^{ij} \sqrt{\det g} \, \de^2 x
  = \frac{1}{8 \pi} \int_C \Phi \Delta \Phi \, \de \vol \, .
\end{equation}
Here $\Delta$ denotes the positive-definite Laplace operator on $C$, acting on functions,
given by
\begin{equation}
  \Delta = \de^* \de = - \frac{1}{\sqrt{\det g}} \partial_i \left(\sqrt{\det g}\,g^{ij}\partial_j\right).
\end{equation}

An important feature here is that the action 
is \ti{conformally invariant}:
for any function $\omega: C \to \bbR$, 
if we rescale the metric on $C$ by
\begin{equation} \label{eq:conformal-rescaling}
  g \to \e^{2 \omega} g \, ,
\end{equation}
the action $S(\Phi)$ for any given $\Phi$ is unchanged.

\begin{exercise}
Verify the conformal invariance of \eqref{eq:free-boson-action}. Would it still be true in dimensions other than $2$?
\end{exercise}

It is useful to study the critical points of $S$.
Making a small variation $\Phi \to \Phi + \delta \Phi$ 
we get
\begin{equation}
  \delta S = \frac{1}{8 \pi} \int_C (\delta \Phi \Delta \Phi + \Phi \Delta (\delta \Phi)) \de \vol = \frac{1}{4 \pi} \int_C \delta \Phi \Delta \Phi \, \de \vol
\end{equation}
This procedure defines the \ti{variational derivative}
\begin{equation}
  \frac{\delta S}{\delta \Phi(x)} = \frac{1}{4 \pi} \Delta \Phi(x) \, .
\end{equation}
Requiring it to vanish
gives the \ti{equation of motion}
\begin{equation} \label{eq:scalar-equation-of-motion}
  \Delta \Phi = 0 \, .
\end{equation}
It is equivalent to requiring that
$\delta S$ vanishes for all (say, smooth) variations
$\delta \Phi$.


\subsection{Complex structure}

A \ti{conformal structure} on $C$ means a Riemannian metric $g$
up to local rescalings \eqref{eq:conformal-rescaling}. 
Now recall that
a conformal structure and orientation on $C$ 
together are equivalent to a complex structure,
making $C$ into a Riemann surface.
The operator of multiplication by $\I$ in $T^* C$ 
is induced by a rotation by a quarter-turn clockwise, which is
also equivalent to the Hodge $\star$ operator
acting on 1-forms.\footnote{Relative 
to an oriented orthonormal basis of $T^* C$ we have
$\star(1) = e_1 \wedge e_2$, $\star(e_1) = e_2$, $\star(e_2) = -e_1$, $\star (e_1 \wedge e_2) = 1$.}

Thus the conformal invariance
of \eqref{eq:free-boson-action} is equivalent to saying that it depends only on 
the complex structure, not on the whole Riemannian metric.
We can rewrite it to make this more manifest:\footnote{$(\Delta f) \de \vol = -2 \I \partial \bar\partial f$.}
\begin{equation}
S = \frac{\I}{4 \pi} \int_C \partial \Phi \wedge \bar\partial \Phi \, .
\end{equation}

\subsection{Holomorphic factorization}
The equation of motion \eqref{eq:scalar-equation-of-motion} says, at least locally, that $\Phi$ is the
real part of a holomorphic function: 
\begin{equation}
  \Phi = \phi + \bphi, \qquad \bar\partial \phi = 0 \, .
\end{equation}

This is a hint that the theory will have some kind of factorization
into holomorphic and antiholomorphic parts.
However, note that $\phi$ is not uniquely determined
by $\Phi$: we have the freedom to shift $\phi$ by an imaginary
constant.
This suggests that the factorization will not be perfect.

% This might lead one to suspect \fixme{... give a good heuristic for why the factorization occurs at level of 
% correlation functions}


\subsection{Currents}

We have two conserved currents: if we write
\begin{equation}
  j = \de \Phi, \qquad \tj = \star \de \Phi
\end{equation}
then using \eqref{eq:scalar-equation-of-motion} we have
\begin{equation}
  \de j = 0, \qquad \de \tj = 0 \, .
\end{equation}

Equivalently, we can 
decompose
\begin{equation}
  j = J + \bJ
\end{equation}
according to type, and the two pieces are separately closed: 
\begin{equation}
\bar\partial J = 0, \qquad \partial \bJ = 0 \, .
\end{equation}


\subsection{The energy-momentum tensor}

We define
\begin{equation}
  T_{ij}(x) = \frac{2}{\sqrt{\det g}} \frac{\delta S}{\delta g^{ij}(x)} \, .
\end{equation}
This is the classical energy-momentum tensor. 
A direct computation gives
\begin{equation} \label{eq:classical-T-free-boson}
 T_{ij} = \frac{1}{4\pi} \left( \partial_i \Phi \partial_j \Phi - \frac12 g_{ij} \partial_k \Phi \partial^k \Phi \right)
\end{equation}

\begin{exercise}
Verify \eqref{eq:classical-T-free-boson}.
\end{exercise}

Again
$T_{ij}$ is a conserved current, in the sense that the equations of motion imply
\begin{equation} \label{eq:T-conserved}
  \nabla^i T_{ij} = 0 \, .
\end{equation}

\begin{exercise}
Verify \eqref{eq:T-conserved} for the energy-momentum tensor
\eqref{eq:classical-T-free-boson}.
\end{exercise}

\begin{exercise}
The equation \eqref{eq:T-conserved} holds in any field theory
which is diffeomorphism invariant. Prove it.
(Hint: consider a vector field
$\xi$ on $C$ giving an infinitesimal diffeomorphism; use the
fact that $S$
is invariant under the simultaneous variation of all 
fields and the metric induced by $\xi$; integrate by parts.)
\end{exercise}

In particular, the trace $T_i^i = 0$, reflecting the 
conformal invariance.




\subsection{Recollection on Gaussian integrals}

Suppose $\cF = \bbR^N$
with coordinates $x^1, \dots, x^N$,
and
\begin{equation} \label{eq:0d-action-multifield}
  S = \half x^i M_{ij} x^j
\end{equation}
where $M$ is symmetric. Then we have
\begin{equation} \label{eq:gaussian-det}
  \DIP{1} = \int_{\cF} \de \vec{x} \, \e^{-\half x^i M_{ij} x^j} = \frac{(2\pi)^{N/2}}{\sqrt{\det M}}.
\end{equation}

\begin{exercise}
Prove that $\DIP{x^{i_1} \cdots x^{i_k}} = 0$ whenever $k$ is odd,
and compute directly 
\begin{equation}
\frac{\DIP{x^i x^j}}{\IP{1}} = (M^{-1})^{ij} \, .  
\end{equation}
\end{exercise}

More generally, $2m$-point functions are given by \fixme{...}


\subsection{Correlation functions of currents}

Now we start discussing correlation functions.
For this the global structure of $C$ matters.
Suppose $C$ is compact connected.

We formally would have
\begin{equation}
  \DIP{1} = \frac{1}{\sqrt{\det \Delta'}}
\end{equation}
where $\Delta'$ is the Laplacian acting on $\cF$
(recall this is the space of functions with integral zero;
equivalently it is the orthocomplement of the constant function,
with respect to the $L^2$ metric).

There is a problem:
the eigenvalues of $\Delta'$
diverge to $\infty$, so $\det \Delta'$ has no chance of 
existing in the ordinary sense.

\begin{exercise} Suppose $C$ is a rectangular torus, with
side lengths $R_1$ and $R_2$. Write out the eigenvalues of 
$\Delta'$. 
\end{exercise}

We postpone dealing with the UV divergence for now, by considering
\ti{normalized} correlation functions, where the divergence will
cancel out. For instance,
let us consider the 2-point correlation function
\begin{equation}
  \frac{\DIP{J(z) J(w)}}{\DIP{1}}
\end{equation}
According to our rules in \fixme{...} this should be given 
by \fixme{...}


\subsection{The \texorpdfstring{$0$}{0}-point function}

The energy-momentum tensor determines the variation of 
correlation functions. In particular we expect formally
\begin{equation}
  \DIP{T_{ij}(x)} = \frac{2}{\sqrt{\det g}} \frac{\delta}{\delta g^{ij}(x)} \DIP{1} \, .
\end{equation}
We use this condition in its logarithmic form
\begin{equation}
  \frac{\DIP{T_{ij}(x)}}{\IP{1}} = \frac{2}{\sqrt{\det g}} \frac{\delta}{\delta g^{ij}(x)} \log \IP{1}
\end{equation}
\fixme{...}



\subsection{TFT and factorization}

Let's consider now a 3-dimensional field theory.
Let $M$ be a compact oriented 3-manifold, now without a
Riemannian metric.
Instead of a scalar $\Phi$ we have a $1$-form $a$, and we write the action
\begin{equation}
S(a) = \int_M a \wedge \de a \, .  
\end{equation}
(We could call this ``Chern-Simons theory with gauge group $\bbR$.'')
$S$ is invariant under $a \to a + \de \chi$ where $\chi$
is any function: indeed the variation is
\begin{equation}
  \delta S = \int_M \de \chi \wedge \de a = \int_M \de(\chi \wedge \de a) = \int_{\partial M} \chi \wedge \de a
\end{equation}
which vanishes for $M$ without boundary.
The equations of motion require $\de a = 0$.
Thus the space of classical solutions modulo gauge symmetry
is $H^1(M,\bbR)$.

Then we let $M = C \times I$, for $I$ an interval.
To formulate the field theory on $M$ we need (elliptic) 
boundary conditions at the two ends of $I$.
One possibility which works is to add an additional
coupling on the boundary, of the form
\begin{equation}
  \pm \int_{\partial M} \norm{a}^2
\end{equation}
Note that this coupling uses a conformal structure on $C$.



Now one of the principles of Lagrangian field theory is that an equation which holds classically
also holds in correlation functions. (You can prove it formally by integration by parts.) 
So in particular we expect $\bar\partial J = 0$, $\partial \bJ = 0$ to hold in correlation 
functions. Thus, for instance, we can compute correlation functions
\begin{equation}
  \DIP{J(z_1) \cdots J(z_k) \bJ(w_1) \cdots \bJ(w_l)}
\end{equation}
and they should be holomorphic 1-forms in all of the variables $z_i$,
antiholomorphic 1-forms in all of the variables $w_i$.
More precisely they are holomorphic \ti{away from} the diagonals where some $z_i = z_j$
or $w_i = w_j$. 

A second key property of these correlation functions is the \ti{operator product expansion}:
this says (in correlation functions) that around $z_1 = z_2$ we have
\begin{equation}
  J(z_1) J(z_2) = \frac{1}{(z_1 - z_2)^2} + \regular
\end{equation}
\fixme{give path integral derivation}
\fixme{this is the OPE of Heisenberg VOA}



\appendix


\section{Conventions on Lie algebras}

For a simple Lie algebra $\fg$,
there is a canonical invariant bilinear form
$\kappa$. We normalize it so that the long roots of $\fg$
have length $2$.

The quadratic Casimiar operator is defined by 
\begin{equation}
  C_2 = \sum_i X_i X_i
\end{equation}
where $(X_i)$ is an orthonormal basis of $\fg$.

The dual Coxeter number $h^\vee$ is defined by the property that
$C_2$ acts in the adjoint representation by the scalar $2 h^\vee$.

When $\fg = \fsl(n)$, the normalized form is
\begin{equation}
  \kappa(X,Y) = \Tr (XY)
\end{equation}
and
\begin{equation}
  h^\vee = n \, .
\end{equation}


% \section{Vertex algebras} \label{sec:vertex-algebras}

% Definition (vertex algebra).  Here \vertex is a complex vector space.

% We informally defined vertex algebras in \fixme{...} above.
% To formalize things, we will have a vector space $\vertex$ of
% defects, and then for $\psi, \chi \in \vertex$
% we will introduce a symbol ``$Y(\psi,z) \chi$'' which should be understood to mean ``the power series expansion 
% of $\psi(z) \chi(0)$ around $z = 0$.''
% Each term in this expansion will be an element in $\vertex$.
% Then we make a lot of computations involving formal Laurent series in many variables,
% but always with the idea that they secretly represent the expansions of actual objects
% which are analytic away from diagonals.
% In that situation what you expect is that there are two different expansions, one in the region
% $\abs{z} > \abs{w}$, another in the region $\abs{w} > \abs{z}$.
% So the symbol I wrote $\psi(z) \phi(w) \chi$ will get represented by \ti{two} different things in this 
% formalization.
% But they are closely related: they are $Y(\psi,z) Y(\phi,w) \chi$ and $Y(\phi, w) Y(\psi,z) \chi$.



\section{The zeta-regularized determinant} \label{sec:zeta-determinant}
 
In the theory of the free boson one expects that
\begin{equation}
  \IP{1} = \sqrt{A} / \sqrt{\det{}' \Delta}
\end{equation}
Since $\Delta$ acts on an infinite-dimensional space we need to explain what we mean by
its determinant.
Precisely we mean the \ti{zeta-regularized} determinant.

Define 
\begin{equation}
 \zeta_\Delta(s) = \sum_{k \ge1 }\lambda_k^{-s} 
\end{equation}
for
$\re(s) > 1$, where $0=\lambda_0<\lambda_1\le\lambda_2\le\cdots$ are the eigenvalues of $\Delta$.
In fact this function admits meromorphic continuation to the full
$s$-plane. To see this consider the operator $\e^{-t \Delta}$ 
giving heat flow for time $t$. This is a nice
integral operator with a smooth kernel $K_t(x,y)$, i.e.
\begin{equation}
  \e^{-t \Delta} f = \int_{y \in C} K_t(x,y) f(y) \, .
\end{equation}
As $t \to 0$ the kernel $K_t$ 
concentrates around the diagonal (it converges in some distributional sense to a $\delta$-function). On the diagonal we have\footnote{To understand the leading term $\frac{1}{4 \pi t}$, note the heat kernel on $\bbR^2$ is a Gaussian, $K_t(x,y) = \frac{1}{4 \pi t} \exp\left(- \frac{\abs{x-y}^2}{4t} \right)$.}
\begin{equation}
  K_t(x,x) \sim \frac{1}{4 \pi t} \left(1 + \frac{R}{6} t + \cdots\right)
\end{equation}
We can subtract these leading divergent terms from $K_t$,
to get a new kernel $\widetilde K_t(x) = K_t(x,x) - \frac{1}{4 \pi t} - \frac{R}{24 \pi}$.
Now, for $\Re(s) > 1$ we have an integral representation
\begin{equation}
  \zeta_\Delta(s) = \frac{1}{\Gamma(s)} \int_0^\infty \frac{\de t}{t} t^s \Tr{}' \e^{-t \Delta}
\end{equation}

With the analytic continuation computed, we finally have
\begin{equation}
\det{}'\Delta = \exp\big(-\zeta'_{\Delta}(0)\big)
\end{equation}




\printbibliography

\end{document}





\section{Free BV theories}

We write the free boson in the BV language.
Very generally this would mean we should write
$S = \frac12 \IP{\phi, Q \phi}$
where $(\phi,Q)$ is an elliptic complex,
and $\IP{,}$ is integration of some local density-valued
pairing of degree $-1$, so that $\IP{,}$ is invariant for $Q$.
Then $Q \phi = 0$ is the equation of motion.
In the case of the free boson we should have fields and antifields,
i.e. $\phi = (\Phi, \Phi^+)$ 
and $Q \Phi = 0$, $Q \Phi^+ = \Delta \Phi$.

Pick $Q^{GF}$ so that $[Q,Q^{GF}] = \Delta$
and propagator
$P = \int_\eps^L (Q^{GF} \otimes 1) K_t \de t$
where $K_t$ is the heat kernel.
(Naively the propagator would be
$\int_0^\infty K_t \de t$, reflecting the identity
$\int_0^\infty \exp(-t \lambda) \de t = \lambda^{-1}$.)



\section{General structure of CFT}

\subsection{The energy-momentum tensor}

We have operators $T_{ij}$ with some universal structure.

The trace has $\IP{T^i_i} = \frac{c}{24 \pi} R$.

The holomorphic part has the transformation law 
\begin{equation} \label{eq:virasoro-coordinate-change}
  T(p)^{z} = \left( \frac{\de w(p)}{\de z(p)} \right)^2 T(p)^{w} + \frac{c}{12} \{w,z\} \, ,
\end{equation}
where $\{\cdot,\cdot\}$ denotes the Schwarzian derivative:
$\{f(z),z\} = \frac{f'''(z)}{f'(z)} - \frac32 \left( \frac{f''(z)}{f'(z)} \right)^2$.






% To add:
% -1/8pi is correct normalization for free boson action if we want J = d\Phi to have 1/(x-y)^2 OPE
% according to \cite{DiFrancescoMathieuSenechal1997CFT}

% Polchinski explains why the trace of T is (-c/12) R [or maybe -c/(24\pi) R, check conventions!]
% [integrated version of this is Polyakov's formula I think]
% calculation: supose we know the trace is a.R for some a
% conservation of T says \nabla^z T_{zz} = -a/2 \partial_z R
% then make Weyl transformations on both sides
% more generally, without assuming anything about the trace in advance,
% this formula should come from computing second Weyl 
% variation of the correlation function in
% two ways: either take the T-T OPE or not

% correlation functions in the full CFT use a local coordinate system around each insertion,
% much as for the correlation functions in the chiral halves

% term "conformal blocks" originates in
% \cite{BPZ1984}



For the 0-point function we know the holomorphic factorization is not quite complete until we expand in conformal blocks
For the 1-point function? Should be zero after averaging
For the 2-point function? e.g. <J Jbar> in genus 1? It's constant I think, something like 1/(im \tau); can calculate it directly using the Green's function, or indirectly using the conformal blocks; the answers are the same
[the un-normalized conformal-block calculation would integrate a^2 exp(-(im\tau) a^2), that gives (im \tau)^(-3/2),
so normalized is (im \tau)^{-1}]




To compute JJ OPE:
free field 





For example, if $C$ is a surface of genus $1$, one has \fixme{fix area-dependent factor}
\begin{equation}
  \det{}'\Delta = (\im \tau) \abs{\eta(\tau)}^4
\end{equation}
where
\begin{equation}
  \eta(\tau) = q^{1/24} \prod_{n=1}^\infty (1 - q^n), \qquad q = e^{2 \pi \I \tau}
\end{equation}

This zeta-regularization might seem ad hoc.
A more physical method of regularization is as follows.
We use the \ti{formal} identity
\begin{equation}
\log \det A = -\int_0^\infty \frac{\de t}{t} \Tr(\e^{-t A})
\end{equation}
This identity is always nonsense, because of a 
divergence near $t = 0$. However it is true ``up to 
an infinite constant,'' in the sense that for $a$, $b$ positive
\begin{equation}
  \log a - \log b = -\int_0^\infty \frac{\de t}{t} (\e^{-t a} - \e^{-t b})
\end{equation}

Then we get a regularized determinant
\begin{equation}
   \log \det{}' \Delta = - \int_0^\infty \frac{\de t}{t} \int_C \widetilde K_t(x)
\end{equation}
\fixme{explain its relation to the zeta-regularized one}

The factorization of this determinant is proved in
\cite{McIntyreTakhtajan2006}.








\subsection{Conformal blocks for affine Kac-Moody algebras}

Continue with the setup of the last section.
Given a chiral correlation system we could consider just the 
sector with $m=0$, i.e. the quantities
\begin{equation}
  \IP{\prod_i \chi_i(p_i)}
\end{equation}
which give a map
\begin{equation}
  \bigotimes_i \vrep_i \to \bbC \, .
\end{equation}
Now suppose $C$ is compact. In that case
this map is not arbitrary: it obeys some constraints,
coming from the existence of the meromorphic 1-forms
\begin{equation}
 \eta = \IP{J^X(\cdot) \prod_i \chi_i(p_i)} \, .
\end{equation}
To formulate the constraints,
let $f$ be any global section of $\cO(*p)$, i.e. 
any meromorphic function on $C$, holomorphic away from 
the $p_i$. Then $f \eta$ is a meromorphic $1$-form
on $C$, and thus
\begin{equation}
  \sum_{i=1}^n \Res_{p_i} (f \eta) = 0 \, .
\end{equation}
So what are these residues?
Laurent expand near each $p_i$ as
\begin{equation}
  f = \sum_{n \in \bbZ} f_{i,n}\, z_i^{n}
\end{equation}
Then, a short computation gives that
\begin{equation}
  \Res_{p_j} f \eta = \sum_{n \in \bbZ} \IP{\left( \prod_{i \neq j} \chi_i(p_i)  \right) f_{j,n} J^X_n \chi_j(p_j)}
\end{equation}
(note only finitely many terms are nonzero).

This motivates the following.
For $X\in\fg$ define a formal infinite sum
\begin{equation}
  X(f)_i = \sum_{n \in \bbZ} f_{i,n}\, J^{X}_{n} ,
\end{equation}
and let $X(f) \in \End(\bigotimes_i \vrep_i)$ be\footnote{The operator $X(f)_i$ depends on the local coordinate $z_i$ we chose; different choices differ by conjugation by an operator in
$\End(\vrep_i)$, 
built by exponentiating a Virasoro action. Thus in the end
the spaces $\Coinv(C,\vertex)$ for different choices of local 
coordinate are canonically isomorphic. \fixme{watch out for choice of log in defining $\lambda^{L_0}$}}\footnote{For any vector $v \in \vrep_i$ we always have $J^X_n v = 0$ for large enough $n$, so $X(f)_i v$ is well defined.}
\begin{equation}
  X(f) = \sum_{i=1}^n \, 
  \bid\otimes\cdots\otimes X(f)_i \otimes\cdots\otimes \bid \, .
\end{equation}

\begin{defn}[Kac-Moody vertex algebra coinvariants on a Riemann surface]
The space of coinvariants is
\begin{equation}
  \Coinv(C, \vertex)
  = \Bigl(\bigotimes_{i=1}^n \vrep_i\Bigr)\Big/\Big\langle
  X(f)\,v \ \Big|\ X\in\fg,\ f\in H^0(\cO(*p)),\
  v\in\bigotimes_i \vrep_i
  \Big\rangle
\end{equation}
\end{defn}

\begin{defn}[Conformal block for Kac-Moody vertex algebra]
\begin{equation}
  \Conf(C, \vertex) = \Coinv(C, \vertex)^* \, .
\end{equation}
In other words, a conformal block is a map
$\bigotimes_i \vrep_i \to \bbC$ which annihilates all
elements of the form $X(f) v$.
\end{defn}

What we just explained is:

\begin{prop}[Chiral correlation systems give conformal blocks]
Given a chiral correlation system for $\vertex = V^k(\fg)$,
the map
\begin{equation}
  \bigotimes_i \vrep_i \to \bbC \, .
\end{equation}
given by
\begin{equation}
  \bigotimes_i \chi_i(p_i) \mapsto  \IP{\prod_i \chi_i(p_i)}
\end{equation}
is a conformal block.
\end{prop}

The converse is also true \fixme{...}

